\documentclass[12pt]{article}
\usepackage[utf8]{inputenc}
\usepackage[T1]{fontenc}
\usepackage{hyperref}
\title{Vishavan Language Resources for K-12 Curriculum Development}
\author{Compiled by Claude Sonnet 5 (claude-sonnet-5) for Andr\'as Kornai}
\date{}
\begin{document}
\maketitle

\section{Introduction}
Vishavan (ISO 639-3 \texttt{vis}; also Malankudi, Malarkuti) is a Dravidian language of a tribal community in central Kerala, India, settled around Idamalyar (Ernakulam District) and Vazhachal (Thrissur District). Wikipedia reports about 150 speakers (1994 figure). Search results showed no educational materials.

\section{Primary Educational Resources}
None found.

\section{Linguistic and Reference Materials}
\begin{itemize}
\item Vishavan language (Wikipedia): \url{https://en.wikipedia.org/wiki/Vishavan_language}
\item Endangered Languages Project entry: \url{https://endangeredlanguages.com/lang/3397}
\item Endangered Languages Project bibliography: \url{https://endangeredlanguages.com/lang/3397/bibliography}
\end{itemize}

\section{Religious and Community Resources}
None found.

\section{Digital and Audio Resources}
None found.

\section{Curriculum Development Context}
Vishavan is not known to be taught in schools and has no found orthography or reference grammar. The bibliography above is the best starting point for locating documentation; its contents were not verified. Malayalam is the likely regional school language.
\end{document}